\section{Discussion}
\markright{Discussion}

The forward-bias readings show the expected nonlinear behavior of the silicon p-n
junction diode, with current increasing steadily as the applied voltage rises and the
depletion barrier is reduced. However, a discrepancy was observed between the
theoretical cut-in voltage of $0.7\,\text{V}$ and the experimental results, where
the diode voltage ($V_F$) only reached a maximum of $0.48\,\text{V}$. This is
attributable to the $10\,\text{k}\Omega$ series resistor being too large, which
limited the forward current to a range where the diode remains in the "knee"
region of its characteristic curve. In reverse bias, the current remained
extremely small; with the $10\,\text{k}\Omega$ resistor, our multimeter did not
display measurable reverse current, which is consistent with the diode blocking
majority carrier flow. If a $1\,\text{k}\Omega$ resistor had been used, the
full $0.7\,\text{V}$ conduction threshold and a larger reverse leakage current
would likely have been more observable. Small differences between simulation
and observation are attributable to real-diode effects such as tolerance, meter
loading, and internal resistance.
